% \documentclass[../main.tex]{subfiles}
% \begin{document}
\section{Background and Context}
\subsection{What Are Algorithms?}
Algorithms are commonly defined as finite, ordered sets of instructions designed to accomplish a specific task or solve a particular problem.
Within computer science, algorithms provide the procedural logic that directs how information is processed, manipulated, and evaluated.
Educationally, algorithms are significant because they introduce students to structured reasoning and formal approaches to problem
  solving while allowing students to practice the necessary syntax required for programming~\cite{grover2017assessing,nijenhuis2023teaching}.

In secondary education, teaching algorithms involves more than instructing students to write code however.
Algorithms represent the underlying logic and sequence of operations that guide computational processes, whereas programming syntax refers to the
  specific grammatical rules of a programming language used to express those ideas~\cite{csta2017}.
Students may therefore understand how to design an algorithm without necessarily being proficient in a language such as Python or Java.
For example, a student can describe the logical process for sorting a list of numbers or determining the largest value in a dataset using plain
  language, diagrams, or pseudocode before implementing the procedure syntactically in
  software~\cite{nijenhuis2023teaching,futschek2006algorithmic}.

The construction of an algorithm therefore requires students to define goals, analyze constraints, test possible solutions, and revise their completed work when errors occur. 
This aligns closely with established educational theories surrounding inquiry-based learning and metacognition~\cite{grover2017assessing}.
Through algorithmic practice, students learn that complex problems can often be approached through iterative refinement, experimentation, and logical
  organization~\cite{grover2017assessing}.
As such, Algorithms often feature prominently in secondary computer science educational standards~\cite{csta2017}.

\subsection{Growth of Secondary Computer Science Education}

Secondary computer science education has expanded significantly over the past two decades as schools, politicians, and employers increasingly
  recognize the importance of computer science as a path of success~\cite{grover2017assessing,csta2017}.
In many secondary contexts, computing instruction has historically been associated with digital literacy and basic software use, while more recent computer science education standards place greater emphasis on computational thinking, algorithms, abstraction, and problem solving~\cite{grover2017assessing,csta2017}.


National and state-level efforts have increasingly integrated computer science into secondary education through revised standards, expanded curricula, and greater attention to teacher preparation~\cite{yadav2016computational,csta2017}.
Many educational systems now promote computer science as a core academic discipline comparable to mathematics, biology or chemistry rather than as a
  specialized elective reserved for a small group of students~\cite{csta2017}.
As a result, secondary schools have increasingly introduced courses involving programming, computational thinking, robotics, cybersecurity, video
  game development and data science~\cite{yadav2016computational}.

\subsection{Challenges in Teaching Algorithms}

Teaching algorithms in secondary education presents several significant challenges related to student perception, prior knowledge, and instructional
  design.
Although algorithms are foundational to computer science, many students encounter difficulties when attempting to understand abstract computational
  processes and translate logical procedures into structured solutions.
These challenges are particularly pronounced among new students who may lack prior exposure to formal reasoning, programming concepts, or
  computational thinking in general~\cite{nijenhuis2023teaching}.
As a result, educators must often balance conceptual instruction with scaffolded support to help students develop confidence and understanding.

One major challenge involves student misconceptions about algorithms and programming.
Many secondary students initially perceive algorithms as synonymous with computer code rather than understanding them as a procedure or set of
  instructions for solving a problem.
This misconception can lead students to focus excessively on memorizing syntax while neglecting understanding of how they are to arrive at a logical
  and correct solution~\cite{gal2004efficiency}.
Students may also incorrectly assume that there is only one correct algorithmic approach to a problem, limiting creativity and flexibility in
  problem-solving processes.
Additional misconceptions frequently emerge around concepts such as sequencing, iteration, variables, scope, syntax and logical
  computation~\cite{futschek2006algorithmic,grover2017assessing}.

Algorithmic thinking often requires students to conceptualize invisible computational processes, mentally simulate execution steps, and anticipate
  outcomes across multiple conditions.
Unlike many traditional classroom subjects that involve observable or concrete phenomena, algorithms frequently operate at a level of abstraction
  unfamiliar to secondary students.
Developmental research also suggests that many secondary students are still refining higher-order reasoning skills associated with abstract and
  formal logic~\cite{fitrianna2025adolescent}.
These difficulties can become more pronounced when instruction advances too quickly from concrete examples to generalized computational
  models~\cite{grover2017assessing}.

To further complicate instruction, the relationship between mathematics and computer science can further complicate education on algorithms.
Many teenagers are wary of mathematics as a subject in general, which may transfer into programming and algorithmic problem solving because both
  disciplines involve symbolic reasoning, procedural thinking, and logical analysis~\cite{khoche12024mathematical}.
Students who have previously struggled with mathematical reasoning may approach algorithmic tasks with reduced confidence or heightened apprehension,
  particularly when algorithms are perceived as requiring advanced mathematical knowledge, even though many algorithmic problems rely more heavily on
  logical reasoning and procedural thinking than on complex mathematical computation~\cite{Voog2012,grover2017assessing}.
Additionally, students sometimes experience transfer issues in which mathematical reasoning skills do not automatically apply to programming
  contexts, for example, a student may understand arithmetic in mathematics but struggle to apply similar logic in code, leading to students to begin perceiving algorithmic concepts with apprehension even though such students were initially confident.

% \end{document}
